% Source: https://artofproblemsolving.com/community/c3h2121119_number_theory_questionhow_many_pairs
% Stable identity: aops-topic-2121119-q01
Only the units digits matter. The units digits of the powers repeat with period $4$:
\[\begin{array}{c|cccc}
\text{exponent}\pmod 4&1&2&3&0\\\hline
3^m&3&9&7&1\\
7^n&7&9&3&1
\end{array}\]
The sum is divisible by $10$ exactly when the units digits add up to $10$: $3+7$, $9+1$, $7+3$ or $1+9$. So, according to $n\bmod 4=1,0,3,2$, the exponent $m$ must satisfy $m\bmod4=1,2,3,0$ respectively. Thus for every $n$ there is exactly one residue class of $m$ modulo $4$ that works.

Among $1,\dots,100$ each residue class modulo $4$ contains exactly $25$ numbers, so each of the $205-101+1=105$ values of $n$ gives $25$ values of $m$. The number of pairs is
\[105\cdot25=\boxed{2625}.\]
